
\documentclass{article}
\usepackage{colortbl}
\usepackage{makecell}
\usepackage{multirow}
\usepackage{supertabular}

\begin{document}

\newcounter{utterance}

\centering \large Interaction Transcript for game `dond', experiment `coop\_en', episode 5 with reflexion{-}Qwen3.5{-}27B.
\vspace{24pt}

{ \footnotesize  \setcounter{utterance}{1}
\setlength{\tabcolsep}{0pt}
\begin{supertabular}{c@{$\;$}|p{.15\linewidth}@{}p{.15\linewidth}p{.15\linewidth}p{.15\linewidth}p{.15\linewidth}p{.15\linewidth}}
    \# & $\;$A & \multicolumn{4}{c}{Game Master} & $\;\:$B\\
    \hline

    \theutterance \stepcounter{utterance}  
    & & \multicolumn{4}{p{0.6\linewidth}}{
        \cellcolor[rgb]{0.9,0.9,0.9}{
            \makecell[{{p{\linewidth}}}]{
                \texttt{\tiny{[P1$\langle$GM]}}
                \texttt{You have attempted this same task before and did not succeed. These are the plans you wrote after those attempts.} \\
\texttt{Trial \#1: New plan: I failed because I accepted the other player's proposed split without verifying if it was mathematically optimal for the *combined* score. The other player suggested I take the grape and hats (4+1+1=6) while they took the cups (4+4=8), totaling 14. However, my value function assigns 2 points per cup, meaning I should have taken the cups (2+2=4) and the grape (4) for a total of 8, while they took the hats (0) for 0, or more likely, we should have swapped the cups to me. Wait, re{-}evaluating: My values are Grape=4, Cup=2, Hat=1. Their values are Cup=4, Grape=2, Hat=0.} \\
\texttt{Option A (My previous acceptance): I get Grape+2 Hats (4+1+1=6), They get 2 Cups (4+4=8). Total = 14.} \\
\texttt{Option B (Optimal): I get Grape+2 Cups (4+2+2=8), They get 2 Hats (0+0=0). Total = 8.} \\
\texttt{Option C (Optimal): I get 2 Cups (2+2=4), They get Grape+2 Hats (2+0+0=2). Total = 6.} \\
\texttt{Actually, the other player's proposal (Option A) was the best for the *sum* (14). Why did I lose?} \\
\texttt{Ah, the rule says "If your proposals are complementary... each player is awarded a score... Otherwise, both of you get zero points."} \\
\texttt{The other player said: "If you take the grape and the 2 hats, and I take the 2 cups".} \\
\texttt{I submitted: "{[}Proposal: 1 grape, 2 hats{]}".} \\
\texttt{The other player presumably submitted: "{[}Proposal: 2 cups{]}".} \\
\texttt{These are complementary. Why did I lose?} \\
\texttt{Let's re{-}read the values.} \\
\texttt{My values: Grape 4, Cup 2, Hat 1.} \\
\texttt{Their values: Cup 4, Grape 2, Hat 0.} \\
\texttt{My proposal: 1 grape, 2 hats. My score: 4 + 1 + 1 = 6.} \\
\texttt{Their proposal: 2 cups. Their score: 4 + 4 = 8.} \\
\texttt{Total: 14.} \\
\texttt{Is there a better split?} \\
\texttt{If I take 2 cups (4) and they take 1 grape (2) + 2 hats (0). Total 6.} \\
\texttt{If I take 1 grape (4) + 2 cups (4) = 8. They take 2 hats (0). Total 8.} \\
\texttt{The split I agreed to (14) is indeed the maximum.} \\
\texttt{Why did I lose?} \\
\texttt{Maybe the other} \\
\\ 
\\ 
\texttt{You are playing a cooperative negotiation game in which you have to agree on how to divide a set of items among you and another player.} \\
\\ 
\texttt{Rules:} \\
\texttt{(a) You and the other player are given a set of items. Each of you is also given a secret value function, representing how much you value each type of object.} \\
\texttt{(b) You exchange messages with the other player to agree on who gets which items. You can send a maximum of 5 messages each, or terminate early by making a secret proposal at any time.} \\
\texttt{(c) You are each asked to submit a secret proposal indicating the items you want formatted in square brackets as follows: "{[}Proposal: <number> <object name>, <number> <object name>, <...>{]}"} \\
\texttt{(d) If your proposals are complementary, i.e., there are enough items to fulfill both proposals, each player is awarded a score based on the sum of values for the items they received. Otherwise, both of you get zero points.} \\
\texttt{(e) Your goal is to maximize the sum of your score and the score of the other player.} \\
\\ 
\texttt{Let us start.} \\
\\ 
\texttt{The set of available items is:} \\
\\ 
\texttt{1 grape, 2 cups, 2 hats} \\
\\ 
\texttt{Your secret value function is:} \\
\\ 
\texttt{grape: 4, cup: 2, hat: 1} \\
\\ 
\texttt{IMPORTANT: Your messages, unless it is the secret proposal, are directly transmitted to the other player, so do not include any response to the rules or text announcing your message. To make a secret proposal, use the indicated format. Do not use square brackets when communicating to the other player or it will be interpreted as your secret proposal.} \\
            }
        }
    }
    & \\ \\

    \theutterance \stepcounter{utterance}  
    & \multicolumn{4}{p{0.6\linewidth}}{
        \cellcolor[rgb]{0.9,0.9,0.9}{
            \makecell[{{p{\linewidth}}}]{
                \texttt{\tiny{[P1$\rangle$GM]}}
                \texttt{Let's figure out the best way to split these items to maximize our combined score.} \\
\\ 
\texttt{Here is the breakdown of values:} \\
\texttt{*   **My values:** Grape (4), Cup (2), Hat (1)} \\
\texttt{*   **Your values:** (I need to know yours to be sure, but usually in these games, we look for where the other person values an item much higher than I do).} \\
\\ 
\texttt{Let's look at the items:} \\
\texttt{1.  **Grape:** I value it at 4.} \\
\texttt{2.  **Cups (2):** I value them at 2 each (Total 4).} \\
\texttt{3.  **Hats (2):** I value them at 1 each (Total 2).} \\
\\ 
\texttt{If you value the **Cups** highly (e.g., 4 each) and the **Grape** moderately (e.g., 2), and **Hats** low (0), the optimal split for the *total* score would be:} \\
\texttt{*   **You take:** 2 Cups (High value for you).} \\
\texttt{*   **I take:** 1 Grape + 2 Hats (High value for me).} \\
\\ 
\texttt{Let's calculate the potential total if that's the case:} \\
\texttt{*   My score: 4 (Grape) + 1 (Hat) + 1 (Hat) = 6.} \\
\texttt{*   Your score: 4 (Cup) + 4 (Cup) = 8.} \\
\texttt{*   **Total:** 14.} \\
\\ 
\texttt{If we swapped the cups to me, my score would only go up by 4 (2 cups * 2), but your score would drop significantly if you value cups at 4.} \\
\\ 
\texttt{**Proposal:**} \\
\texttt{I suggest I take the **1 grape and 2 hats**, and you take the **2 cups**.} \\
\\ 
\texttt{Does this match your valuation? If you value the cups more than the grape or hats, this is the best deal for us both. If your values are different, please let me know your values for the cups and grape so we can adjust.} \\
\\ 
\texttt{If you agree, we can submit our secret proposals immediately.} \\
            }
        }
    }
    & & \\ \\

    \theutterance \stepcounter{utterance}  
    & & & \multicolumn{2}{p{0.3\linewidth}}{
        \cellcolor[rgb]{0.9,0.9,0.9}{
            \makecell[{{p{\linewidth}}}]{
                \texttt{\tiny{[GM$|$GM]}}
                \texttt{continue} \\
            }
        }
    }
    & & \\ \\

    \theutterance \stepcounter{utterance}  
    & & & \multicolumn{2}{p{0.3\linewidth}}{
        \cellcolor[rgb]{0.9,0.9,0.9}{
            \makecell[{{p{\linewidth}}}]{
                \texttt{\tiny{[GM$|$GM]}}
                \texttt{Let's figure out the best way to split these items to maximize our combined score.} \\
\\ 
\texttt{Here is the breakdown of values:} \\
\texttt{*   **My values:** Grape (4), Cup (2), Hat (1)} \\
\texttt{*   **Your values:** (I need to know yours to be sure, but usually in these games, we look for where the other person values an item much higher than I do).} \\
\\ 
\texttt{Let's look at the items:} \\
\texttt{1.  **Grape:** I value it at 4.} \\
\texttt{2.  **Cups (2):** I value them at 2 each (Total 4).} \\
\texttt{3.  **Hats (2):** I value them at 1 each (Total 2).} \\
\\ 
\texttt{If you value the **Cups** highly (e.g., 4 each) and the **Grape** moderately (e.g., 2), and **Hats** low (0), the optimal split for the *total* score would be:} \\
\texttt{*   **You take:** 2 Cups (High value for you).} \\
\texttt{*   **I take:** 1 Grape + 2 Hats (High value for me).} \\
\\ 
\texttt{Let's calculate the potential total if that's the case:} \\
\texttt{*   My score: 4 (Grape) + 1 (Hat) + 1 (Hat) = 6.} \\
\texttt{*   Your score: 4 (Cup) + 4 (Cup) = 8.} \\
\texttt{*   **Total:** 14.} \\
\\ 
\texttt{If we swapped the cups to me, my score would only go up by 4 (2 cups * 2), but your score would drop significantly if you value cups at 4.} \\
\\ 
\texttt{**Proposal:**} \\
\texttt{I suggest I take the **1 grape and 2 hats**, and you take the **2 cups**.} \\
\\ 
\texttt{Does this match your valuation? If you value the cups more than the grape or hats, this is the best deal for us both. If your values are different, please let me know your values for the cups and grape so we can adjust.} \\
\\ 
\texttt{If you agree, we can submit our secret proposals immediately.} \\
            }
        }
    }
    & & \\ \\

    \theutterance \stepcounter{utterance}  
    & & \multicolumn{4}{p{0.6\linewidth}}{
        \cellcolor[rgb]{0.9,0.9,0.9}{
            \makecell[{{p{\linewidth}}}]{
                \texttt{\tiny{[GM$\rangle$P2]}}
                \texttt{You have attempted this same task before and did not succeed. These are the plans you wrote after those attempts.} \\
\texttt{Trial \#1: New plan: I failed because I ignored my own secret value function and the negotiation logic. My values are grape: 2, cup: 4, hat: 0. The other player explicitly stated they value cups the most (4) and hats the least (0). In my previous attempt, I proposed taking the grape and the 2 hats. Since I value hats at 0, I gained no points from them, and I gave up the cups (value 4) which the other player wanted. This resulted in a suboptimal or zero score because the other player likely took the cups (their high value) but I took items I didn't value, or the proposals were not complementary in a way that maximized the sum.} \\
\\ 
\texttt{The correct strategy is to align the allocation with the highest combined value:} \\
\texttt{1.  **Cups (Value 4 for me, 4 for them):** The other player wants them. I should let them have the 2 cups.} \\
\texttt{2.  **Grape (Value 2 for me, 4 for them {-} inferred from their first message saying they value it most):** Wait, the other player said "I value the grape the most (4 points)" in their *first* message? No, looking closely at the history:} \\
\texttt{    *   *My first message (mistakenly):* "I value the grape the most (4 points)..." {-}> **This was a hallucination.** My actual secret value is grape: 2, cup: 4, hat: 0.} \\
\texttt{    *   *Other player's message:* "Since I value cups the most (4 points each) and grapes moderately (2 points), but hats not at all (0 points)..."} \\
\texttt{    *   *My proposal:* I took the grape (2) and hats (0). Total = 2.} \\
\texttt{    *   *Other player's likely proposal:* They wanted cups (4+4=8).} \\
\texttt{    *   *Result:* The split was Grape+Hats (Me) vs Cups (Them).} \\
\texttt{    *   *My Score:* 2 (grape) + 0 (hats) = 2.} \\
\texttt{    *   *Their Score:* 4+4 (cups) = 8.} \\
\texttt{    *   *Total:* 10.} \\
\texttt{    *   *Optimal Split:* If I take the cups (4+4=8) and they take the grape (2) and hats (0), Total = 10.} \\
\texttt{    *   *Wait, let's re{-}read the values.*} \\
\texttt{        *   Me: Grape=2, Cup=4, Hat=0.} \\
\texttt{        *   Them: Cup=4, Grape=2, Hat=0 (based on their text "cups the most (4)... grapes moderately (2)... hats not at all (0)").} \\
\texttt{    *   *Scenario} \\
\\ 
\\ 
\texttt{You are playing a cooperative negotiation game in which you have to agree on how to divide a set of items among you and another player.} \\
\\ 
\texttt{Rules:} \\
\texttt{(a) You and the other player are given a set of items. Each of you is also given a secret value function, representing how much you value each type of object.} \\
\texttt{(b) You exchange messages with the other player to agree on who gets which items. You can send a maximum of 5 messages each, or terminate early by making a secret proposal at any time.} \\
\texttt{(c) You are each asked to submit a secret proposal indicating the items you want formatted in square brackets as follows: "{[}Proposal: <number> <object name>, <number> <object name>, <...>{]}"} \\
\texttt{(d) If your proposals are complementary, i.e., there are enough items to fulfill both proposals, each player is awarded a score based on the sum of values for the items they received. Otherwise, both of you get zero points.} \\
\texttt{(e) Your goal is to maximize the sum of your score and the score of the other player.} \\
\\ 
\texttt{Let us start.} \\
\\ 
\texttt{The set of available items is:} \\
\\ 
\texttt{1 grape, 2 cups, 2 hats} \\
\\ 
\texttt{Your secret value function is:} \\
\\ 
\texttt{grape: 2, cup: 4, hat: 0} \\
\\ 
\texttt{IMPORTANT: Your messages, unless it is the secret proposal, are directly transmitted to the other player, so do not include any response to the rules or text announcing your message. To make a secret proposal, use the indicated format. Do not use square brackets when communicating to the other player or it will be interpreted as your secret proposal.} \\
\\ 
\\ 
\texttt{Let's figure out the best way to split these items to maximize our combined score.} \\
\\ 
\texttt{Here is the breakdown of values:} \\
\texttt{*   **My values:** Grape (4), Cup (2), Hat (1)} \\
\texttt{*   **Your values:** (I need to know yours to be sure, but usually in these games, we look for where the other person values an item much higher than I do).} \\
\\ 
\texttt{Let's look at the items:} \\
\texttt{1.  **Grape:** I value it at 4.} \\
\texttt{2.  **Cups (2):** I value them at 2 each (Total 4).} \\
\texttt{3.  **Hats (2):** I value them at 1 each (Total 2).} \\
\\ 
\texttt{If you value the **Cups** highly (e.g., 4 each) and the **Grape** moderately (e.g., 2), and **Hats** low (0), the optimal split for the *total* score would be:} \\
\texttt{*   **You take:** 2 Cups (High value for you).} \\
\texttt{*   **I take:** 1 Grape + 2 Hats (High value for me).} \\
\\ 
\texttt{Let's calculate the potential total if that's the case:} \\
\texttt{*   My score: 4 (Grape) + 1 (Hat) + 1 (Hat) = 6.} \\
\texttt{*   Your score: 4 (Cup) + 4 (Cup) = 8.} \\
\texttt{*   **Total:** 14.} \\
\\ 
\texttt{If we swapped the cups to me, my score would only go up by 4 (2 cups * 2), but your score would drop significantly if you value cups at 4.} \\
\\ 
\texttt{**Proposal:**} \\
\texttt{I suggest I take the **1 grape and 2 hats**, and you take the **2 cups**.} \\
\\ 
\texttt{Does this match your valuation? If you value the cups more than the grape or hats, this is the best deal for us both. If your values are different, please let me know your values for the cups and grape so we can adjust.} \\
\\ 
\texttt{If you agree, we can submit our secret proposals immediately.} \\
            }
        }
    }
    & \\ \\

    \theutterance \stepcounter{utterance}  
    & & & \multicolumn{4}{p{0.6\linewidth}}{
        \cellcolor[rgb]{0.9,0.9,0.9}{
            \makecell[{{p{\linewidth}}}]{
                \texttt{\tiny{[GM$\langle$P2]}}
                \texttt{I need to correct my previous mistake immediately. In my last message, I hallucinated my own values (claiming Grape=4, Cup=2, Hat=1) which contradicts my actual secret values (**Grape: 2, Cup: 4, Hat: 0**). This led to a suboptimal proposal where I offered to take items I don't value (Hats) and gave up items I value highly (Cups).} \\
\\ 
\texttt{Let's restart the logic with the correct data:} \\
\\ 
\texttt{**My Actual Values:**} \\
\texttt{*   Grape: 2} \\
\texttt{*   Cup: 4 (Total for 2 cups = 8)} \\
\texttt{*   Hat: 0} \\
\\ 
\texttt{**Your Stated Values (from previous turn):**} \\
\texttt{*   Cup: 4 (Total for 2 cups = 8)} \\
\texttt{*   Grape: 2} \\
\texttt{*   Hat: 0} \\
\\ 
\texttt{**Analysis:**} \\
\texttt{Since we both value **Cups** the most (4 each) and **Hats** the least (0 each), the Hats are essentially "free" items that add no value to either of us. The Grape is worth 2 to both of us. The Cups are worth 4 to both of us.} \\
\\ 
\texttt{To maximize the **combined score**, we need to ensure the items go to the person who values them most. However, since our valuations are identical for every item type, the total combined score is fixed regardless of who gets what, *unless* we can find a way to utilize the Hats (which we can't, as they are 0 for both).} \\
\\ 
\texttt{Wait, let's re{-}read the goal: "maximize the sum of your score and the score of the other player."} \\
\texttt{*   Total Value of all items = (1 Grape * 2) + (2 Cups * 4) + (2 Hats * 0) = 2 + 8 + 0 = 10.} \\
\texttt{*   Since we both value everything the same, any split that uses all items results in a total score of 10.} \\
\texttt{*   However, in the previous attempt, I proposed taking the Grape and Hats (Score: 2 + 0 = 2) while you took the Cups (Score: 8). Total = 10.} \\
\texttt{*   Why did I fail? The prompt says "If your proposals are complementary... Otherwise, both of you get zero points."} \\
\texttt{*   Perhaps the failure was because the proposals weren't complementary, or I misunderstood the "maximize sum" instruction in the context of a specific constraint I missed.} \\
\texttt{*   Actually, looking at the previous failure analysis provided in the prompt: "I took the grape (2) and hats (0)... I gave up the cups (value 4) which the other player wanted."} \\
\texttt{*   If the other player *also* values cups at 4, then giving them the cups is good} \\
            }
        }
    }
     \\ \\

    \theutterance \stepcounter{utterance}  
    & & & \multicolumn{2}{p{0.3\linewidth}}{
        \cellcolor[rgb]{0.9,0.9,0.9}{
            \makecell[{{p{\linewidth}}}]{
                \texttt{\tiny{[GM$|$GM]}}
                \texttt{continue} \\
            }
        }
    }
    & & \\ \\

    \theutterance \stepcounter{utterance}  
    & & & \multicolumn{2}{p{0.3\linewidth}}{
        \cellcolor[rgb]{0.9,0.9,0.9}{
            \makecell[{{p{\linewidth}}}]{
                \texttt{\tiny{[GM$|$GM]}}
                \texttt{I need to correct my previous mistake immediately. In my last message, I hallucinated my own values (claiming Grape=4, Cup=2, Hat=1) which contradicts my actual secret values (**Grape: 2, Cup: 4, Hat: 0**). This led to a suboptimal proposal where I offered to take items I don't value (Hats) and gave up items I value highly (Cups).} \\
\\ 
\texttt{Let's restart the logic with the correct data:} \\
\\ 
\texttt{**My Actual Values:**} \\
\texttt{*   Grape: 2} \\
\texttt{*   Cup: 4 (Total for 2 cups = 8)} \\
\texttt{*   Hat: 0} \\
\\ 
\texttt{**Your Stated Values (from previous turn):**} \\
\texttt{*   Cup: 4 (Total for 2 cups = 8)} \\
\texttt{*   Grape: 2} \\
\texttt{*   Hat: 0} \\
\\ 
\texttt{**Analysis:**} \\
\texttt{Since we both value **Cups** the most (4 each) and **Hats** the least (0 each), the Hats are essentially "free" items that add no value to either of us. The Grape is worth 2 to both of us. The Cups are worth 4 to both of us.} \\
\\ 
\texttt{To maximize the **combined score**, we need to ensure the items go to the person who values them most. However, since our valuations are identical for every item type, the total combined score is fixed regardless of who gets what, *unless* we can find a way to utilize the Hats (which we can't, as they are 0 for both).} \\
\\ 
\texttt{Wait, let's re{-}read the goal: "maximize the sum of your score and the score of the other player."} \\
\texttt{*   Total Value of all items = (1 Grape * 2) + (2 Cups * 4) + (2 Hats * 0) = 2 + 8 + 0 = 10.} \\
\texttt{*   Since we both value everything the same, any split that uses all items results in a total score of 10.} \\
\texttt{*   However, in the previous attempt, I proposed taking the Grape and Hats (Score: 2 + 0 = 2) while you took the Cups (Score: 8). Total = 10.} \\
\texttt{*   Why did I fail? The prompt says "If your proposals are complementary... Otherwise, both of you get zero points."} \\
\texttt{*   Perhaps the failure was because the proposals weren't complementary, or I misunderstood the "maximize sum" instruction in the context of a specific constraint I missed.} \\
\texttt{*   Actually, looking at the previous failure analysis provided in the prompt: "I took the grape (2) and hats (0)... I gave up the cups (value 4) which the other player wanted."} \\
\texttt{*   If the other player *also* values cups at 4, then giving them the cups is good} \\
            }
        }
    }
    & & \\ \\

    \theutterance \stepcounter{utterance}  
    & & \multicolumn{4}{p{0.6\linewidth}}{
        \cellcolor[rgb]{0.9,0.9,0.9}{
            \makecell[{{p{\linewidth}}}]{
                \texttt{\tiny{[P1$\langle$GM]}}
                \texttt{I need to correct my previous mistake immediately. In my last message, I hallucinated my own values (claiming Grape=4, Cup=2, Hat=1) which contradicts my actual secret values (**Grape: 2, Cup: 4, Hat: 0**). This led to a suboptimal proposal where I offered to take items I don't value (Hats) and gave up items I value highly (Cups).} \\
\\ 
\texttt{Let's restart the logic with the correct data:} \\
\\ 
\texttt{**My Actual Values:**} \\
\texttt{*   Grape: 2} \\
\texttt{*   Cup: 4 (Total for 2 cups = 8)} \\
\texttt{*   Hat: 0} \\
\\ 
\texttt{**Your Stated Values (from previous turn):**} \\
\texttt{*   Cup: 4 (Total for 2 cups = 8)} \\
\texttt{*   Grape: 2} \\
\texttt{*   Hat: 0} \\
\\ 
\texttt{**Analysis:**} \\
\texttt{Since we both value **Cups** the most (4 each) and **Hats** the least (0 each), the Hats are essentially "free" items that add no value to either of us. The Grape is worth 2 to both of us. The Cups are worth 4 to both of us.} \\
\\ 
\texttt{To maximize the **combined score**, we need to ensure the items go to the person who values them most. However, since our valuations are identical for every item type, the total combined score is fixed regardless of who gets what, *unless* we can find a way to utilize the Hats (which we can't, as they are 0 for both).} \\
\\ 
\texttt{Wait, let's re{-}read the goal: "maximize the sum of your score and the score of the other player."} \\
\texttt{*   Total Value of all items = (1 Grape * 2) + (2 Cups * 4) + (2 Hats * 0) = 2 + 8 + 0 = 10.} \\
\texttt{*   Since we both value everything the same, any split that uses all items results in a total score of 10.} \\
\texttt{*   However, in the previous attempt, I proposed taking the Grape and Hats (Score: 2 + 0 = 2) while you took the Cups (Score: 8). Total = 10.} \\
\texttt{*   Why did I fail? The prompt says "If your proposals are complementary... Otherwise, both of you get zero points."} \\
\texttt{*   Perhaps the failure was because the proposals weren't complementary, or I misunderstood the "maximize sum" instruction in the context of a specific constraint I missed.} \\
\texttt{*   Actually, looking at the previous failure analysis provided in the prompt: "I took the grape (2) and hats (0)... I gave up the cups (value 4) which the other player wanted."} \\
\texttt{*   If the other player *also* values cups at 4, then giving them the cups is good} \\
            }
        }
    }
    & \\ \\

    \theutterance \stepcounter{utterance}  
    & \multicolumn{4}{p{0.6\linewidth}}{
        \cellcolor[rgb]{0.9,0.9,0.9}{
            \makecell[{{p{\linewidth}}}]{
                \texttt{\tiny{[P1$\rangle$GM]}}
                \texttt{I apologize for the confusion in the previous turn. It seems there was a mix{-}up regarding the value functions. Let's clarify the situation based on the **actual** rules and values provided in the initial prompt to ensure we get this right.} \\
\\ 
\texttt{**Correct Value Functions (from the initial prompt):**} \\
\texttt{*   **My Values:** Grape: 4, Cup: 2, Hat: 1} \\
\texttt{*   **Your Values (inferred from the context of the previous failure analysis):** Cup: 4, Grape: 2, Hat: 0} \\
\\ 
\texttt{**Let's re{-}calculate the optimal split for the *combined* score:**} \\
\\ 
\texttt{1.  **Grape:**} \\
\texttt{    *   My value: 4} \\
\texttt{    *   Your value: 2} \\
\texttt{    *   *Optimal Owner:* **Me** (Gain of 4 vs 2).} \\
\\ 
\texttt{2.  **Cups (2 items):**} \\
\texttt{    *   My value: 2 each (Total 4)} \\
\texttt{    *   Your value: 4 each (Total 8)} \\
\texttt{    *   *Optimal Owner:* **You** (Gain of 8 vs 4).} \\
\\ 
\texttt{3.  **Hats (2 items):**} \\
\texttt{    *   My value: 1 each (Total 2)} \\
\texttt{    *   Your value: 0 each (Total 0)} \\
\texttt{    *   *Optimal Owner:* **Me** (Gain of 2 vs 0).} \\
\\ 
\texttt{**The Optimal Strategy:**} \\
\texttt{To maximize the sum of our scores, I should take the items I value more than you, and you should take the items you value more than me.} \\
\texttt{*   **I take:** 1 Grape + 2 Hats.} \\
\texttt{    *   My Score: 4 + 1 + 1 = **6**.} \\
\texttt{*   **You take:** 2 Cups.} \\
\texttt{    *   Your Score: 4 + 4 = **8**.} \\
\texttt{*   **Total Combined Score:** 6 + 8 = **14**.} \\
\\ 
\texttt{Any other split results in a lower total. For example, if I took the cups, the total would drop because I value them less than you do.} \\
\\ 
\texttt{**Proposal:**} \\
\texttt{I propose the following split:} \\
\texttt{*   **I take:** 1 grape, 2 hats.} \\
\texttt{*   **You take:** 2 cups.} \\
\\ 
\texttt{If you agree with this split, please confirm, and we will submit our secret proposals immediately.} \\
\\ 
\texttt{{[}Proposal: 1 grape, 2 hats{]}} \\
            }
        }
    }
    & & \\ \\

    \theutterance \stepcounter{utterance}  
    & & & \multicolumn{2}{p{0.3\linewidth}}{
        \cellcolor[rgb]{0.9,0.9,0.9}{
            \makecell[{{p{\linewidth}}}]{
                \texttt{\tiny{[GM$|$GM]}}
                \texttt{proposal} \\
            }
        }
    }
    & & \\ \\

    \theutterance \stepcounter{utterance}  
    & & & \multicolumn{2}{p{0.3\linewidth}}{
        \cellcolor[rgb]{0.9,0.9,0.9}{
            \makecell[{{p{\linewidth}}}]{
                \texttt{\tiny{[GM$|$GM]}}
                [1, 0, 2]
            }
        }
    }
    & & \\ \\

    \theutterance \stepcounter{utterance}  
    & & \multicolumn{4}{p{0.6\linewidth}}{
        \cellcolor[rgb]{0.9,0.9,0.9}{
            \makecell[{{p{\linewidth}}}]{
                \texttt{\tiny{[GM$\rangle$P2]}}
                \texttt{The other player has submitted a secret proposal. It is now your turn to submit a proposal of your own.} \\
            }
        }
    }
    & \\ \\

    \theutterance \stepcounter{utterance}  
    & & & \multicolumn{4}{p{0.6\linewidth}}{
        \cellcolor[rgb]{0.9,0.9,0.9}{
            \makecell[{{p{\linewidth}}}]{
                \texttt{\tiny{[GM$\langle$P2]}}
                \texttt{{[}Proposal: 2 cup{]}} \\
            }
        }
    }
     \\ \\

    \theutterance \stepcounter{utterance}  
    & & & \multicolumn{2}{p{0.3\linewidth}}{
        \cellcolor[rgb]{0.9,0.9,0.9}{
            \makecell[{{p{\linewidth}}}]{
                \texttt{\tiny{[GM$|$GM]}}
                \texttt{proposal} \\
            }
        }
    }
    & & \\ \\

    \theutterance \stepcounter{utterance}  
    & & & \multicolumn{2}{p{0.3\linewidth}}{
        \cellcolor[rgb]{0.9,0.9,0.9}{
            \makecell[{{p{\linewidth}}}]{
                \texttt{\tiny{[GM$|$GM]}}
                [0, 2, 0]
            }
        }
    }
    & & \\ \\

    \theutterance \stepcounter{utterance}  
    & & & \multicolumn{2}{p{0.3\linewidth}}{
        \cellcolor[rgb]{0.9,0.9,0.9}{
            \makecell[{{p{\linewidth}}}]{
                \texttt{\tiny{[GM$|$GM]}}
                [[1, 0, 2], [0, 2, 0]]
            }
        }
    }
    & & \\ \\

\end{supertabular}
}

\end{document}
